\chapter{Curvilinear coordinate transformation}\label{app:CoordinateTransform}
A curvilinear coordinate transformation is employed to enable the calculation of non-rectangular bodies of water.
Two grids are introduced: a curvilinear $xy$-grid, following the curvature of the shallow water body and a computational/numerical $\xi \eta$-grid, created by mapping the curvilinear $xy$-grid to an orthogonal coordinate system via a coordinate transformation.
The global Cartesian $xy$-coordinate system is used for reference of the numerical grid.
An example of this is sketched in \autoref{fig:grids}.
%
\begin{figure}[H]
    \begin{center}
        \def\svgwidth{0.80\textwidth} % scaling text
        \resizebox{0.99\textwidth}{!}{
            \input{figures/coordinate_transformation.pdf_tex}
        }
    \end{center}
	\caption{Mesh mapping from the physical Cartesian $xy$-grid to the numerical $\xi\eta$-grid and vice versa.}
\label{fig:grids}
\end{figure}
%
Following the chain rule and assuming time-independent grids, we can express the differential operators in the global coordinate system as functions of the differential operators in the body-fitted curvilinear grid as follows:
%
\begin{align}\label{1ordertransx}
	\pdiff{}{x} &= \pdiff{\xi}{x} \pdiff{}{\xi}+\pdiff{\eta}{x}\pdiff{}{\eta},
\\
 \label{1ordertransy}
	\pdiff{}{y} &= \pdiff{\xi}{y}\pdiff{}{\xi}+\pdiff{\eta}{y}\pdiff{}{\eta},
\end{align}
in matrix-vector notation these equations reads:
\begin{align}
    \begin{pmatrix} \pdiff{}{x} \\ \pdiff{}{y} \end{pmatrix}
    =
    \begin{pmatrix}
        \pdiff{\xi}{x} &
        \pdiff{\eta}{x} \\
        \pdiff{\xi}{y} &
        \pdiff{\eta}{y}
    \end{pmatrix}
    \begin{pmatrix} \pdiff{}{\xi} \\ \pdiff{}{\eta} \end{pmatrix}.
    \label{eq:grid_transformation}
\end{align}
%
where the transformation coefficients $\pdiff{\xi}{x}$, $\pdiff{\eta}{x}$, $\pdiff{\xi}{y}$, $\pdiff{\eta}{y}$ are unknown and but can be determined by the grid definitions.

The goal is to express the transformation coefficients of \autoref{eq:grid_transformation} in terms of its inverted derivatives (grid definitions).
The inverse transformation reads:
%
\begin{align}\label{}
    \begin{pmatrix} \pdiff{}{\xi} \\ \pdiff{}{\eta} \end{pmatrix}
    =
    \begin{pmatrix}
        \pdiff{x}{\xi} &
        \pdiff{y}{\xi} \\
        \pdiff{x}{\eta} &
        \pdiff{y}{\eta}
    \end{pmatrix}
    \begin{pmatrix} \pdiff{}{x} \\ \pdiff{}{y} \end{pmatrix}.
    \label{eq:inverse_grid_transformation}
\end{align}
%
We introduce a shorthand notation for the transformation coefficients with a subscript, e.g.\ $\xi_x = \pdiff{\xi}{x}$ and calculate the inverse of the right-hand side to give:
%
\begin{align}
	\begin{pmatrix}
		x_\xi &
		y_\xi \\
		x_\eta &
		y_\eta
	\end{pmatrix}^{-1}
	=
	\frac{1}{x_{\xi}y_{\eta}-y_{\xi}x_{\eta}}
	\begin{pmatrix}
		y_{\eta} &
		-y_{\xi}\\
		-x_{\eta} &
		x_{\xi}
	\end{pmatrix}
	=
	\frac{1}{J}
	\begin{pmatrix}
		y_{\eta} &
		-y_{\xi}\\
		-x_{\eta} &
		x_{\xi}
	\end{pmatrix}
    \label{eq:inverse_transformation}
\end{align}
with the Jacobian determinant defined as:
\begin{align}
J =
\begin{vmatrix}
    \pdiff{x}{\xi} & \pdiff{x}{\eta} \\
    \pdiff{y}{\xi} & \pdiff{y}{\eta}
\end{vmatrix}
= x_{\xi}y_{\eta}-y_{\xi}x_{\eta}.
\end{align}
leading to
\begin{align}
	\xi_x=\frac{1}{J}y_{\eta}, \quad \eta_x = -\frac{1}{J}y_{\xi}, \quad
	\xi_y=-\frac{1}{J}x_{\eta}, \quad \eta_y = \frac{1}{J}x_{\xi}.
    \label{eq:inverse_transformation_1}
\end{align}
where $J = x_{\xi}y_{\eta}-x_{\eta}y_{\xi}$ is the determinant of the coordinate transformation matrix.

Hence, we can use \autoref{eq:inverse_transformation_1} to transform the partial derivatives (\autoref{eq:grid_transformation}) leading to:
	\begin{align}\label{trans}
		\pdiff{}{x}
		&=\xi_x\pdiff{}{\xi}
		+ \eta_x\pdiff{}{\eta}
		=\frac{1}{J}\left(y_{\eta} \pdiff{}{\xi}
		-y_{\xi}\pdiff{}{\eta}\right)
\\
		\pdiff{}{y}
		&=\xi_y\pdiff{}{\xi}
		+\eta_y \pdiff{}{\eta}
		=\frac{1}{J}\left(-x_{\eta}\pdiff{}{\xi}
		+x_{\xi}\pdiff{}{\eta}
		\right)
	\end{align}
\subparagraph*{One dimension}
\begin{quote}
{{\scriptsize
In one dimension we have $y_\xi = x_\eta = 0$, $\xi_x = y_\eta/J$ and $J = x_\xi y_\eta$, thus
\begin{align}
    \pdiff{}{x}
    &=\xi_x\pdiff{}{\xi}
    =\frac{1}{J}y_{\eta} \pdiff{}{\xi}
    = \frac{1}{x_\xi} \pdiff{}{\xi}
\end{align}
}}
\end{quote}
%--------------------------------------------------------------------------------
\paragraph*{Second derivatives}
Second derivatives are instead transformed by applying this operator twice.
This yields:
%
\begin{align}
	\pdiff[2]{}{x}
	&=\frac{1}{J}\Big[y_{\eta} \pdiff{}{\xi} \Big( \frac{1}{J}\Big(y_{\eta} \pdiff{}{\xi}
	-y_{\xi}\pdiff{}{\eta}\Big)\Big)
	-y_{\xi}\pdiff{}{\eta} \Big( \frac{1}{J}\Big(y_{\eta} \pdiff{}{\xi}
	-y_{\xi}\pdiff{}{\eta}\Big) \Big) \Big],
\\
\pdiff[2]{}{y}
		&=\frac{1}{J}\Big[-x_{\eta}\pdiff{}{\xi} \Big( \frac{1}{J}\Big(-x_{\eta}\pdiff{}{\xi}
		+x_{\xi}\pdiff{}{\eta}
		\Big)\Big)
		+x_{\xi}\pdiff{}{\eta} \Big(\frac{1}{J}\Big(-x_{\eta}\pdiff{}{\xi}
		+x_{\xi}\pdiff{}{\eta}
		\Big) \Big)
		\Big],
\end{align}
%
which can be rewritten to:
%
\begin{align}
	\begin{aligned}
		\pdiff[2]{}{x} & =
		\frac{y_{\eta}}{J^2} y_{\eta} \pdiff[2]{}{\xi}
		+ y_{\xi} \frac{y_{\xi}}{J^2} \pdiff[2]{}{\eta}
		- 2\frac{y_{\eta}}{J^2} y_{\xi}\pdiff{^2}{\xi \partial \eta} + \\
		& +  \Big(
		\frac{y_{\xi} y_{\eta}J_{\eta}}{J^3}
		- \frac{y_{\eta}^2 J_{\xi}}{J^3}
		+ \frac{y_{\eta}y_{\xi\eta}}{J^2}
		-  \frac{y_{\eta \eta} y_{\xi}}{J^2}
		\Big)\pdiff{}{\xi} +\\
		& + \Big(
		\frac{y_{\eta} y_{\xi} J_{\xi}}{J^3}
		- \frac{y_{\xi}^2 J_{\eta}}{J^3}
		+  \frac{y_{\xi \eta}y_{\xi}}{J^2}
		- \frac{ y_{\xi\xi} y_{\eta}}{J^2}
		\Big)\pdiff{}{\eta}
	\end{aligned}
\end{align}
\begin{align}
    \begin{aligned}
		\pdiff[2]{}{y} & =
	    x_{\eta}\frac{x_{\eta}}{J^2}\pdiff[2]{}{\xi}
		+ x_{\xi}\frac{x_{\xi}}{J^2}\pdiff[2]{}{\eta}
		- 2 x_{\xi}\frac{x_{\eta}}{J^2}\pdiff{^2}{\xi \partial\eta} + \\
		& + \Big(
		\frac{x_{\xi}x_{\eta}J_{\eta}}{J^3}
		-\frac{x_{\eta}^2J_{\xi}}{J^3}
		+ \frac{x_{\xi\eta} x_{\eta}}{J^2}
		- \frac{x_{\eta\eta}x_{\xi}}{J^2}
		\Big) \pdiff{}{\xi} + \\
		&+ \Big(
		\frac{x_{\eta}x_{\xi}J_{\xi}}{J^3}
		- \frac{ x_{\xi}^2J_{\eta}}{J^3}
		+ \frac{x_{\xi\eta}x_{\xi}}{J^2}
		- \frac{x_{\xi \xi}x_{\eta}}{J^2}
		\Big) \pdiff{}{\eta},
    \end{aligned}
\end{align}
%
after repeated application of the chain rule, which should equal the application of the chain rule directly:
%
\begin{align}\label{2orderpartialtransxfinal}
	\pdiff[2]{}{x}
	= \xi_x^2 \pdiff[2]{}{\xi}
	+ \eta_x^2 \pdiff[2]{}{\eta}
	+ 2\eta_x\xi_x \pdiff{^2}{\xi \partial \eta}
	+ \xi_{xx} \pdiff{}{\xi}
	+ \eta_{xx} \pdiff{}{\eta},
\end{align}
%
\begin{align}\label{2orderpartialtransyfinal}
	\pdiff[2]{}{y}
	= \xi_y^2 \pdiff[2]{}{\xi}
	+ \eta_y^2 \pdiff[2]{}{\eta}
	+ 2\eta_y\xi_y\pdiff{^2}{\xi \partial \eta}
	+ \xi_{yy} \pdiff{}{\xi}
	+ \eta_{yy}\pdiff{}{\eta}.
\end{align}
%
from which we can find the following relations:
%
\begin{align}
		\xi_{xx} &=
		\frac{y_{\xi} y_{\eta}J_{\eta}}{J^3}
		- \frac{y_{\eta}^2 J_{\xi}}{J^3}
		+ \frac{y_{\eta}y_{\xi\eta}}{J^2}
		-  \frac{y_{\eta \eta} y_{\xi}}{J^2} \\
		\eta_{xx} &=
		\frac{y_{\eta} y_{\xi} J_{\xi}}{J^3}
		- \frac{y_{\xi}^2 J_{\eta}}{J^3}
		+  \frac{y_{\xi \eta}y_{\xi}}{J^2}
		- \frac{ y_{\xi\xi} y_{\eta}}{J^2} \\
		\xi_{yy} &=
		\frac{x_{\xi}x_{\eta}J_{\eta}}{J^3}
		-\frac{x_{\eta}^2J_{\xi}}{J^3}
		+ \frac{x_{\xi\eta} x_{\eta}}{J^2}
		- \frac{x_{\eta\eta}x_{\xi}}{J^2} \\
		\eta_{yy} &=
		\frac{x_{\eta}x_{\xi}J_{\xi}}{J^3}
		- \frac{ x_{\xi}^2J_{\eta}}{J^3}
		+ \frac{x_{\xi\eta}x_{\xi}}{J^2}
		- \frac{x_{\xi \xi}x_{\eta}}{J^2}
\end{align}
\subparagraph*{One dimension}
\begin{quote}
{\scriptsize
In one dimension we have $y_\xi = x_\eta = 0$, $\xi_x = y_\eta/J$ and $J = x_\xi y_\eta$, thus
\begin{align}
    \pdiff[2]{}{x}
    & =\xi_x\pdiff{}{\xi}\left(  \xi_x\pdiff{}{\xi} \right)
      =\frac{1}{x_\xi}\pdiff{}{\xi}\left( \frac{1}{x_\xi} \pdiff{}{\xi} \right) =
    \\
    & =\frac{1}{x_\xi^2} \pdiff[2]{}{\xi} - \frac{x_{\xi\xi}}{x_\xi^3}\pdiff{}{\xi}
\end{align}
Multiplying this equation by $\Dx = \Dxi x_\xi^2$ it yields \citet[eq.\ 2]{Borsboom1998}, which reads:
\begin{align}
    \Dx \pdiff[2]{}{x}
& = \Dxi x_\xi^2 \left( \frac{1}{x_\xi^2} \pdiff[2]{}{\xi} - \frac{x_{\xi\xi}}{x_\xi^3}\pdiff{}{\xi} \right)
=
\Dxi \left( \pdiff[2]{}{\xi} - \frac{x_{\xi\xi}}{x_\xi}\pdiff{}{\xi} \right)
\end{align}
}
\end{quote}
%-------------------------------------------------------------------------------
\section{Transforming the terms}\label{sec:transforming_the_terms}
For completeness, we show the transformation of various terms here since some particular choices will be made to ease the implementation.
For example, transforming the spatial derivatives in the continuity equation in a conservative form
%
\begin{align}
	\pdiff{q}{x} + \pdiff{r}{y}  & = \frac{1}{J} \Bigg(  y_{\eta} \pdiff{q}{\xi}
    -y_{\xi}\pdiff{q}{\eta} +\left( - x_{\eta}\pdiff{r}{\xi}
    +x_{\xi}\pdiff{r}{\eta} \right) \Bigg)
\end{align}
and adding $(y_{\eta\xi} - y_{\xi\eta})q = 0$ and $(x_{\eta\xi} - x_{\xi\eta})r = 0$, yields
\begin{align}
    \pdiff{q}{x} + \pdiff{r}{y}
    & = \frac{1}{J} \Bigg(
    y_{\eta} \pdiff{q}{\xi}- y_{\xi}\pdiff{q}{\eta} + {\color{blue}y_{\eta\xi}q - y_{\xi\eta}q}
    \\
    & \qquad - x_{\eta}\pdiff{r}{\xi} + x_{\xi}\pdiff{r}{\eta} +{\color{blue}x_{\eta\xi}r - x_{\xi\eta}r} \Bigg)
    \\
    & =\frac{1}{J}\Bigg( \pdiff{}{\xi} \left( y_{\eta} q - x_{\eta} r \right) + \pdiff{}{\eta} \left( -y_{\xi} q + x_{\xi} r \right)\Bigg).
    \label{eq:full_curvilinear_mass_flux}
\end{align}
%
Taking the finite volume approach, yields
\begin{align}
&\frac{1}{J}\int_{\Omega_{\xi\eta}}
\begin{pmatrix}
    \pdiff{}{\xi} \\ \pdiff{}{\eta}
\end{pmatrix}
\dotp
\begin{pmatrix}
    y_{\eta} q - x_{\eta} r \\ -y_{\xi} q + x_{\xi} r
\end{pmatrix}
\, dl
=
\frac{1}{J}\oint_{\partial\Omega_{\xi\eta}}
\begin{pmatrix}
    y_{\eta} q - x_{\eta} r \\ -y_{\xi} q + x_{\xi} r
\end{pmatrix}
\dotp
\begin{pmatrix}
    \nxi \\ \neta
\end{pmatrix}
\, dl =
\\
&
\frac{1}{J}\oint_{\partial\Omega_{\xi\eta}} \left( \left( y_{\eta} q - x_{\eta} r \right)\nxi + \left( -y_{\xi} q + x_{\xi} r \right) \neta \right)\, dl
\label{eq:transformed_mass_flux}
\end{align}
%
%--------------------------------------------------------------------------------
\subsection{Convection term}
The convection term in the $q$-momentum equation can be transformed as:
%
\begin{align}
	&\pdiff{(q^2/h)}{x} + \pdiff{(qr/h)}{y} =
\\
	&= \frac{1}{J} \left(y_{\eta} \pdiff{(q^2/h)}{\xi} - y_{\xi} \pdiff{(q^2/h)}{\eta} -x_{\eta} \pdiff{(qr/h)}{\xi} + x_{\xi} \pdiff{(qr/h)}{\eta}\right) =
    \label{eq:xi_convection_term_1}
\end{align}
\begin{align}
	\begin{aligned}
		&= \frac{1}{J} \Bigg(y_{\eta} \pdiff{(q^2/h)}{\xi} - y_{\xi} \pdiff{(q^2/h)}{\eta} + {\color{blue}\frac{q^2}{h}  y_{\xi\eta} - \frac{q^2}{h} y_{\eta\xi}} +  \\
		&\quad -x_{\eta} \pdiff{(qr/h)}{\xi} + x_{\xi} \pdiff{(qr/h)}{\eta} {\color{blue} - \frac{qr}{h}  x_{\xi\eta} + \frac{qr}{h} x_{\eta\xi}} \Bigg) =
	\end{aligned}
\end{align}
\begin{align}
	= \frac{1}{J} \pdiff{}{\xi}\left(y_{\eta} \frac{q^2}{h} - x_{\eta} \frac{qr}{h} \right) + \frac{1}{J} \pdiff{}{\eta}\left(-y_{\xi} \frac{q^2}{h} + x_{\xi} \frac{qr}{h} \right).
    \label{eq:xi_convection_term}
\end{align}
%
Taking the finite volume approach and applying Green's theorem, yields
\begin{align}
    \frac{1}{J}\oint_{\partial\Omega_{\xi\eta}} \left( \left( y_{\eta}\frac{q^2}{h} - x_{\eta} \frac{qr}{h} \right) \nxi + \left(-y_{\xi} \frac{q^2}{h} + x_{\xi} \frac{qr}{h} \right) \neta \right)\, dl.
\end{align}
which is equal to (and look to the similarity with \autoref{eq:transformed_mass_flux})
\begin{align}
    \frac{1}{J}\oint_{\partial\Omega_{\xi\eta}} \left[  \left( y_{\eta}q - x_{\eta} r \right) \frac{q}{h}  \nxi + \left( -y_{\xi}q  + x_{\xi}r\right) \frac{q}{h} \neta \right] \, dl.\label{eq:transformed_convection_x}
\end{align}

A similar transformation is valid for the convection term in the $r$-momentum equation, which reads
\begin{align}
    &\pdiff{(rq/h)}{x} + \pdiff{(r^2/h)}{y} =
\\
    &= \frac{1}{J} \left(y_{\eta} \pdiff{(rq/h)}{\xi} - y_{\xi} \pdiff{(rq/h)}{\eta} -x_{\eta} \pdiff{(r^2/h)}{\xi} + x_{\xi} \pdiff{(r^2/h)}{\eta}\right) =
    \label{eq:eta_convection_term_1}
\end{align}
\begin{align}
    \begin{aligned}
        &= \frac{1}{J} \Bigg(y_{\eta} \pdiff{(rq/h)}{\xi} - y_{\xi} \pdiff{(rq/h)}{\eta} + {\color{blue}\frac{rq}{h}  y_{\xi\eta} - \frac{rq}{h} y_{\eta\xi}} +  \\
        &\quad -x_{\eta} \pdiff{(r^2/h)}{\xi} + x_{\xi} \pdiff{(r^2/h)}{\eta} {\color{blue} - \frac{r^2}{h}  x_{\xi\eta} + \frac{r^2}{h} x_{\eta\xi}} \Bigg) =
    \end{aligned}
\end{align}
\begin{align}
    = \frac{1}{J} \pdiff{}{\xi}\left(y_{\eta} \frac{rq}{h} - x_{\eta} \frac{r^2}{h} \right) + \frac{1}{J} \pdiff{}{\eta}\left(-y_{\xi} \frac{rq}{h} + x_{\xi} \frac{r^2}{h} \right).
    \label{eq:eta_convection_term}
\end{align}
Taking the finite volume approach and applying Green's theorem, yields
\begin{align}
    \frac{1}{J}\oint_{\partial\Omega_{\xi\eta}} \left(
    \left(y_{\eta} \frac{rq}{h} - x_{\eta} \frac{r^2}{h} \right) \nxi
    + \left(-y_{\xi} \frac{rq}{h} + x_{\xi} \frac{r^2}{h} \right) \neta
     \right)
    \, dl
\end{align}
which is equal to (and look to the similarity with \autoref{eq:transformed_mass_flux})
\begin{align}
    \frac{1}{J}\oint_{\partial\Omega_{\xi\eta}} \left[
    \left(y_{\eta} q - x_{\eta} r \right)\frac{r}{h} \nxi
    + \left(-y_{\xi} q + x_{\xi} r \right)\frac{r}{h} \neta
    \right]
    \, dl
\end{align}
%\newpage
%--------------------------------------------------------------------------------
\subsection{Viscosity term} \label{sec:viscosity_term_curvilinear}
The viscosity term in vector notation reads (not yet integrated over a control volume):
\begin{align}
    &-\nabla \dotp \left( \nu h \left(  \nabla (\vec{q}/h) + \nabla (\vec{q}^T/h) \right) \right)
    %    =
    %    \\
    %    & \qquad \qquad =
    %    - \oint_{\partial\Omega}  \left( \nu h \left( \nabla (\vec{q}/h) + \nabla (\vec{q}^T/h) \right) \right) \dotp \vecn\, dl
\end{align}
Written in components ($\vec{q} = (q, r)^T$):
\begin{align}
    &
    - % \int_{\Omega}
    \begin{pmatrix}
        \pdiff{}{x} \\
        \pdiff{}{y}
    \end{pmatrix}
    \dotp
    \Bigg[
    \nu h
    \left(
    \begin{pmatrix}
        \pdiff{(q/h)}{x} & \pdiff{(q/h)}{y}
        \\
        \pdiff{(r/h)}{x} & \pdiff{(r/h)}{y}
    \end{pmatrix}
    +
    \begin{pmatrix}
        \pdiff{(q/h)}{x} & \pdiff{(r/h)}{x}
        \\
        \pdiff{(q/h)}{y} & \pdiff{(r/h)}{y}
    \end{pmatrix}
    \right) \Bigg]
    \\
%\, d\omega
    &
    - % \int_{\Omega}
\begin{pmatrix}
    \pdiff{}{x} \\
    \pdiff{}{y}
\end{pmatrix}
\dotp
\Bigg[
\nu h
\begin{pmatrix}
    \pdiff{(q/h)}{x} + \pdiff{(q/h)}{x} & \pdiff{(q/h)}{y} + \pdiff{(r/h)}{x}
    \\
    \pdiff{(r/h)}{x} + \pdiff{(q/h)}{y} & \pdiff{(r/h)}{y} + \pdiff{(r/h)}{y}
\end{pmatrix}
\Bigg]
%\, d\omega
\end{align}
When the viscous value is not symmetric due to the used numerical method
the viscosity tensor is written as:
%
\begin{align}
    \mat{\nu} =not symmetric due to the used numerical method
    the viscosity tensor is written as
    \begin{pmatrix}
        \nu_{11} & \nu_{12} \\
        \nu_{21}  & \nu_{22}
    \end{pmatrix}=
    \begin{pmatrix}
        \nu + \Psi_{11} & \nu + \Psi_{12} \\
        \nu + \Psi_{21}  & \nu + \Psi_{22}
    \end{pmatrix}
\end{align}
with $\Psi_{i,j}, i,j \in \{1,2\}$ representing the added artificial viscosity due to regularization.
Substituting this tensor yields:
\begin{align}
- % \int_{\Omega}
\begin{pmatrix}
    \pdiff{}{x} \\
    \pdiff{}{y}
\end{pmatrix}
\dotp
\Bigg[
h
\begin{pmatrix}
    \nu_{11} & \nu_{12} \\
    \nu_{21}  & \nu_{22}
\end{pmatrix}
\begin{pmatrix}
    2\pdiff{(q/h)}{x} & \pdiff{(q/h)}{y} + \pdiff{(r/h)}{x}
    \\
    \pdiff{(r/h)}{x} + \pdiff{(q/h)}{y} & 2\pdiff{(r/h)}{y}
\end{pmatrix}
\Bigg]
% \, d\omega
\end{align}
%
\begin{align}
- % \int_{\Omega}
    \begin{pmatrix}
        \pdiff{}{x} \\
        \pdiff{}{y}
    \end{pmatrix}
    \dotp
    \left[
    h
 \left(
\begin{array}{ll}
    \nu_{11} \left( 2\pdiff{(q/h)}{x} \right) + &
    \nu_{11} \left( \pdiff{(q/h)}{y} + \pdiff{(r/h)}{x} \right) +
    \\
    \qquad + \nu_{12} \left( \pdiff{(r/h)}{x} + \pdiff{(q/h)}{y} \right) & \qquad + \nu_{12} \left( 2 \pdiff{(r/h)}{y} \right) \\
    \nu_{21} \left( 2\pdiff{(q/h)}{x} \right)  + &
    \nu_{21} \left( \pdiff{(q/h)}{y} + \pdiff{(r/h)}{x} \right) +
    \\
    \qquad + \nu_{22} \left( \pdiff{(r/h)}{x} + \pdiff{(q/h)}{y} \right)
    &
    \qquad + \nu_{22} \left( 2 \pdiff{(r/h)}{y} \right)
\end{array}
\right)
    \right]
\end{align}
\begin{align}
%    \, d\omega
\end{align}
Applying the gradient to the elements of the matrix yields:
%
\newline\textit{First Component:}
\begin{align}
    &-\pdiff{}{x} \left[  2 \nu_{11} h \pdiff{(q/h)}{x} + \nu_{12} h \left(\pdiff{(r/h)}{x} + \pdiff{(q/h)}{y}\right) \right] +
    \\
    &\qquad
    - \pdiff{}{y} \left[ \nu_{11} h \left( \pdiff{(r/h)}{x} + \pdiff{(q/h)}{y}  \right) + 2 \nu_{12} h \pdiff{(r/h)}{y} \right]
\end{align}
\textit{Second Component:}
\begin{align}
    &-\pdiff{}{x} \left[ 2 \nu_{21} h \pdiff{(q/h)}{x} + \nu_{22}  h \left(\pdiff{(r/h)}{x} + \pdiff{(q/h)}{y}\right) \right] +
    \\
    &\qquad
    - \pdiff{}{y} \left[ \nu_{21} h \left(\pdiff{(r/h)}{x} + \pdiff{(q/h)}{y}\right) + 2 \nu_{22} h \pdiff{(r/h)}{y} \right]
\end{align}
%--------------------------------------------------------------------------------
\paragraph*{Viscosity $q$-momentum}
The viscous terms for the $q$-equation read
\begin{align}
    &-\pdiff{}{x} \left( 2 \nu_{11} h \pdiff{(q/h)}{x}
    + \nu_{12} h \left(\pdiff{(r/h)}{x} + \pdiff{(q/h)}{y}\right) \right) +
    \\
    &\qquad
    - \pdiff{}{y} \left( \nu_{11} h \left( \pdiff{(r/h)}{x} + \pdiff{(q/h)}{y} \right) + 2 \nu_{12} h \pdiff{(r/h)}{y} \right)
\end{align}
can be transformed as follows: \newline
\textit{step 1}
\begin{align}
=
- \frac{y_{\eta}}{J} \pdiff{}{\xi} & \left( 2\nu_{11} h\pdiff{(q/h)}{x} \right)
+ \frac{y_{\xi}}{J} \pdiff{}{\eta}\left(2\nu_{11} h\pdiff{(q/h)}{x} \right) +
\\
- \frac{y_{\eta}}{J} \pdiff{}{\xi} & \left( \nu_{12} h \left(\pdiff{(r/h)}{x} + \pdiff{(q/h)}{y}\right) \right)
+ \frac{y_{\xi}}{J} \pdiff{}{\eta} \left( \nu_{12} h \left(\pdiff{(r/h)}{x} + \pdiff{(q/h)}{y}\right) \right) +
\\
+ \frac{x_{\eta}}{J}\pdiff{}{\xi} & \left( \nu_{11} h \left( \pdiff{(r/h)}{x} + \pdiff{(q/h)}{y} \right)\right)
- \frac{x_{\xi}}{J} \pdiff{}{\eta}\left( \nu_{11} h \left( \pdiff{(r/h)}{x} + \pdiff{(q/h)}{y} \right)\right)
\\
+ \frac{x_{\eta}}{J}\pdiff{}{\xi} & \left( 2 \nu_{12} h \pdiff{(r/h)}{y} \right)
- \frac{x_{\xi}}{J} \pdiff{}{\eta}\left( 2 \nu_{12} h \pdiff{(r/h)}{y} \right)
=
\end{align}
Using the chain rule $uv' = (uv)' - u'v$, and assuming $x_{\xi\eta} - x_{\eta\xi} = 0$ and $y_{\xi\eta} - y_{\eta\xi} = 0$, so the terms with second derivative of the coordinate cancel, yields:
%
\begin{align}
    =
     \frac{1}{J} \pdiff{}{\xi} & \Bigg[
     -2\nu_{11} y_{\eta} h\pdiff{(q/h)}{x}
    - \nu_{12} y_{\eta} h \left(\pdiff{(r/h)}{x} + \pdiff{(q/h)}{y}\right)
    \\
    &
    + \nu_{11} x_\eta h \left( \pdiff{(r/h)}{x} + \pdiff{(q/h)}{y} \right)
    + 2 \nu_{12} x_{\eta} h \pdiff{(r/h)}{y}
    \Bigg] +
%    \\
%    &+ \Ccancel[blue]{2 y_{\eta_\xi} \nu_{11} h \pdiff{(q/h)}{x}}
%    - \Cbcancel[blue]{ x_{\eta\xi} \nu_{11} h \left(\pdiff{(q/h)}{y} + \pdiff{(r/h)}{x}\right)}
    \\
    + \frac{1}{J} \pdiff{}{\eta} &\Bigg[ 2\nu_{11} y_{\xi} h \pdiff{(q/h)}{x}
    + \nu_{12} y_{\xi} h \left(\pdiff{(r/h)}{x} + \pdiff{(q/h)}{y}\right) +
    \\
    &
    -  \nu_{11} x_\xi h \left(\pdiff{(r/h)}{x} + \pdiff{(q/h)}{y} \right)
    -  2 \nu_{12} x_{\xi} h \pdiff{(r/h)}{y}
    \Bigg]
%    \\
%    &
%    - \Ccancel[blue]{2 y_{\xi\eta} \nu_{11} h \pdiff{(q/h)}{x}}
%    + \Cbcancel[blue]{x_{\xi\eta} \nu_{11} h \left(\pdiff{(q/h)}{y} + \pdiff{(r/h)}{x}\right)}
\end{align}
\textit{step 2}
\begin{align}
\frac{1}{J} \pdiff{}{\xi}\Bigg[ &
\frac{\nu_{11} h}{J} \left(
- 2 y_{\eta}  y_{\eta} \pdiff{(q/h)}{\xi}
+ 2 y_{\eta} y_{\xi} \pdiff{(q/h)}{\eta} \right)
+
\\ + &
\frac{\nu_{12} h}{J} \left(
- y_{\eta}y_{\eta} \pdiff{(r/h)}{\xi}
+ y_{\eta}y_{\xi} \pdiff{(r/h)}{\eta}
+ y_{\eta}x_{\eta}\pdiff{(q/h)}{\xi}
- y_{\eta}x_{\xi}\pdiff{(q/h)}{\eta} \right) +
\\ + &
\frac{\nu_{11} h}{J} \left(
+ x_{\eta}y_{\eta} \pdiff{(r/h)}{\xi}
- x_{\eta}y_{\xi} \pdiff{(r/h)}{\eta}
- x_{\eta}x_{\eta}\pdiff{(q/h)}{\xi}
+ x_{\eta}x_{\xi}\pdiff{(q/h)}{\eta}
\right) +
\\
+ &
\frac{\nu_{12} h}{J}   \left(
- 2 x_{\eta}x_{\eta} \pdiff{(r/h)}{\xi}
+ 2 x_{\eta}x_{\xi} \pdiff{(r/h)}{\eta}
\right)
\Bigg] +
%---------------
\\
+ \frac{1}{J} \pdiff{}{\eta}\Bigg[ &
\frac{\nu_{11} h}{J} \left(
2 y_{\xi}  y_{\eta} \pdiff{(q/h)}{\xi}
- 2 y_{\xi} y_{\xi} \pdiff{(q/h)}{\eta}
\right)
\\* + &
\frac{\nu_{12} h}{J} \left(
  y_{\xi}y_{\eta} \pdiff{(r/h)}{\xi}
- y_{\xi}y_{\xi} \pdiff{(r/h)}{\eta}
- y_{\xi}x_{\eta}\pdiff{(q/h)}{\xi}
+ y_{\xi}x_{\xi}\pdiff{(q/h)}{\eta}
\right)
\\* + &
\frac{\nu_{11} h}{J} \left(
- x_{\xi}y_{\eta} \pdiff{(r/h)}{\xi}
+ x_{\xi}y_{\xi} \pdiff{(r/h)}{\eta}
+ x_{\xi}x_{\eta}\pdiff{(q/h)}{\xi}
- x_{\xi}x_{\xi}\pdiff{(q/h)}{\eta}
\right)
\\* + &
\frac{\nu_{12} h}{J} \left(
  2 x_{\xi}x_{\eta} \pdiff{(r/h)}{\xi}
- 2 x_{\xi}x_{\xi} \pdiff{(r/h)}{\eta}
\right)
\Bigg]
\label{eq:xi_viscous_term}
\end{align}
%
After applying Green’s theorem with
$\vecn = (\nxi, \neta)^T$ the outward normal vector, and multiplied with $J$  the viscosity term of then $q$-momentum equation reads:
%
\begin{align}
    \Bigg[ &  \frac{\nu_{11}}{J} \left(
    - 2 y_{\eta} y_{\eta} h \pdiff{(q/h)}{\xi}
    + 2 y_{\eta} y_{\xi} h \pdiff{(q/h)}{\eta} \right)
    +
    \\* & +
    \frac{\nu_{12}}{J} \left(
    - y_{\eta}y_{\eta} h \pdiff{(r/h)}{\xi}
    + y_{\eta}y_{\xi} h \pdiff{(r/h)}{\eta}
    + y_{\eta}x_{\eta} h \pdiff{(q/h)}{\xi}
    - y_{\eta}x_{\xi} h \pdiff{(q/h)}{\eta} \right) +
    \\* & +
    \frac{\nu_{11}}{J} \left(
    + x_{\eta}y_{\eta} h \pdiff{(r/h)}{\xi}
    - x_{\eta}y_{\xi} h \pdiff{(r/h)}{\eta}
    - x_{\eta}x_{\eta} h \pdiff{(q/h)}{\xi}
    + x_{\eta}x_{\xi} h \pdiff{(q/h)}{\eta}
    \right) +
    \\* & +
    \frac{\nu_{12}}{J} \left(
    - 2 x_{\eta}x_{\eta} h \pdiff{(r/h)}{\xi}
    + 2 x_{\eta}x_{\xi} h \pdiff{(r/h)}{\eta}
    \right)
    \Bigg] \nxi +
    %---------------
    \\
    + \Bigg[ &
    \frac{\nu_{11}}{J} \left(
    2 y_{\xi}  y_{\eta} h \pdiff{(q/h)}{\xi}
    - 2 y_{\xi} y_{\xi} h \pdiff{(q/h)}{\eta}
    \right)
    \\* & +
    \frac{\nu_{12}}{J} \left(
    y_{\xi}y_{\eta} h \pdiff{(r/h)}{\xi}
    - y_{\xi}y_{\xi} h \pdiff{(r/h)}{\eta}
    - y_{\xi}x_{\eta} h \pdiff{(q/h)}{\xi}
    + y_{\xi}x_{\xi} h \pdiff{(q/h)}{\eta}
    \right)
    \\* & +
    \frac{\nu_{11}}{J} \left(
    - x_{\xi}y_{\eta} h \pdiff{(r/h)}{\xi}
    + x_{\xi}y_{\xi} h \pdiff{(r/h)}{\eta}
    + x_{\xi}x_{\eta} h \pdiff{(q/h)}{\xi}
    - x_{\xi}x_{\xi} h \pdiff{(q/h)}{\eta}
    \right)
    \\* & +
    \frac{\nu_{12}}{J} \left(
    2 x_{\xi}x_{\eta} h \pdiff{(r/h)}{\xi}
    - 2 x_{\xi}x_{\xi} h \pdiff{(r/h)}{\eta}
    \right)
    \Bigg] \neta \label{eq:xi_greens_viscous_term}
\end{align}
%
The linear functions in the first derivatives yields:
\begin{align}
    h \pdiff{(q/h)}{\xi} = \pdiff{q}{\xi} - \frac{q}{h}\pdiff{h}{\xi},
    & \qquad h \pdiff{q/h)}{\eta} = \pdiff{q}{\eta} - \frac{q}{h}\pdiff{h}{\eta},
    \\*
    h \pdiff{(r/h)}{\xi} = \pdiff{r}{\xi} - \frac{r}{h}\pdiff{h}{\xi},
    & \qquad h\pdiff{(r/h)}{\eta} = \pdiff{r}{\eta} - \frac{r}{h}\pdiff{h}{\eta}.
\end{align}
and will be used for the discretization of these viscosity terms.
%--------------------------------------------------------------------------------
\paragraph*{Viscosity $r$-momentum}
The $r$-momentum equation the viscous terms can be transformed as:
\begin{align}
    &-\pdiff{}{x} \left[ 2 \nu_{21} h \pdiff{(q/h)}{x} + \nu_{22}  h \left(\pdiff{(r/h)}{x} + \pdiff{(q/h)}{y}\right) \right] +
    \\
    &\qquad
    - \pdiff{}{y} \left[ \nu_{21} h \left(\pdiff{(r/h)}{x} + \pdiff{(q/h)}{y}\right) + 2 \nu_{22} h \pdiff{(r/h)}{y} \right]
\end{align}
\textit{step 1}
\begin{align}
=&
-\frac{y_{\eta}}{J}\pdiff{}{\xi}\left( 2 \nu_{21} h \pdiff{(q/h)}{x} \right)
+ \frac{y_{\xi}}{J} \pdiff{}{\eta}\left( 2 \nu_{21} h \pdiff{(q/h)}{x}\right)
+
\\
&-
\frac{y_{\eta}}{J}\pdiff{}{\xi}\left( \nu_{22}  h \left(\pdiff{(r/h)}{x} + \pdiff{(q/h)}{y}\right) \right)
+ \frac{y_{\xi}}{J} \pdiff{}{\eta}\left( \nu_{22}  h \left(\pdiff{(r/h)}{x} + \pdiff{(q/h)}{y}\right) \right)
+
\\
&+ \frac{x_{\eta}}{J}\pdiff{}{\xi}\left( \nu_{21} h \left(\pdiff{(r/h)}{x} + \pdiff{(q/h)}{y}\right) \right)
- \frac{x_{\xi}}{J} \pdiff{}{\eta}\left( \nu_{21} h \left(\pdiff{(r/h)}{x} + \pdiff{(q/h)}{y}\right) \right)
+
\\
&+ \frac{x_{\eta}}{J}\pdiff{}{\xi}\left( 2 \nu_{22} h \pdiff{(r/h)}{y} \right)
- \frac{x_{\xi}}{J} \pdiff{}{\eta}\left( 2 \nu_{22} h \pdiff{(r/h)}{y} \right)
\end{align}
Using the chain rule $uv' = (uv)' - u'v$, and assuming $x_{\xi\eta} - x_{\eta\xi} = 0$ and $y_{\xi\eta} - y_{\eta\xi} = 0$, so the terms with second derivative of the coordinate cancel, yields:
\begin{align}
    =
    \frac{1}{J} \pdiff{}{\xi} & \Bigg[
    -2\nu_{21} y_{\eta} h\pdiff{(q/h)}{x}
    - \nu_{22} y_{\eta} h \left(\pdiff{(r/h)}{x} + \pdiff{(q/h)}{y}\right)
    +
    \\*
    &
    + \nu_{21} x_\eta h \left(\pdiff{(r/h)}{x} + \pdiff{(q/h)}{y} \right)
    + 2 \nu_{22} x_{\eta} h \pdiff{(r/h)}{y}
    \Bigg] +
    \\
    + \frac{1}{J} \pdiff{}{\eta} & \Bigg[ 2\nu_{21} y_{\xi} h \pdiff{(q/h)}{x}
    + \nu_{22} y_{\xi} h \left(\pdiff{(r/h)}{x} + \pdiff{(q/h)}{y}\right) +
    \\*
    &
    -  \nu_{21} x_\xi h \left(\pdiff{(r/h)}{x} + \pdiff{(q/h)}{y}\right)
    -  2 \nu_{22} x_{\xi} h \pdiff{(r/h)}{y}
    \Bigg]
\end{align}
%
\textit{step 2}
\begin{align}
= \frac{1}{J} \pdiff{}{\xi} \Bigg[ &
\frac{\nu_{21} h}{J} \left(
- 2 y_{\eta} y_{\eta} h\pdiff{(q/h)}{\xi}
+ 2 y_{\eta} y_{\xi} h\pdiff{(q/h)}{\eta}
\right)
+
\\
+ &\frac{\nu_{22} h}{J} \left(
- y_{\eta} y_{\eta} h \pdiff{(r/h)}{\xi}
+ y_{\eta} y_{\xi} h \pdiff{(r/h)}{\eta}
+ y_{\eta} x_{\eta} h \pdiff{(q/h)}{\xi}
- y_{\eta} x_{\xi} h \pdiff{(q/h)}{\eta}
\right)
+
\\
+ &\frac{\nu_{21} h}{J} \left(
+ x_{\xi} y_{\eta} h \pdiff{(r/h)}{\xi}
- x_{\xi} y_{\xi} h \pdiff{(r/h)}{\eta}
- x_{\xi} x_{\eta} h \pdiff{(q/h)}{\xi}
+ x_{\xi} x_{\xi} h \pdiff{(q/h)}{\eta}
\right)
+
\\
+ &\frac{\nu_{22} h}{J} \left(
- 2 x_{\eta} x_{\eta} h\pdiff{(q/h)}{\xi}
+ 2 x_{\eta} x_{\xi} h\pdiff{(q/h)}{\eta}
\right)
\Bigg] +
%---------------
\\
+ \frac{1}{J} \pdiff{}{\eta}\Bigg[ &
\frac{\nu_{21} h}{J} \left(
  2 y_{\xi} y_{\eta} \pdiff{(q/h)}{\xi}
- 2 y_{\xi} y_{\xi} \pdiff{(q/h)}{\eta}
\right)
\\* + &
\frac{\nu_{22} h}{J} \left(
  y_{\xi}y_{\eta} \pdiff{(r/h)}{\xi}
- y_{\xi}y_{\xi} \pdiff{(r/h)}{\eta}
- y_{\xi}x_{\eta} \pdiff{(q/h)}{\xi}
+ y_{\xi}x_{\xi} \pdiff{(q/h)}{\eta}
\right)
\\* + &
\frac{\nu_{21} h}{J} \left(
- x_{\xi}y_{\eta} \pdiff{(r/h)}{\xi}
+ x_{\xi}y_{\xi} \pdiff{(r/h)}{\eta}
+ x_{\xi}x_{\eta} \pdiff{(q/h)}{\xi}
- x_{\xi}x_{\xi} \pdiff{(q/h)}{\eta}
\right)
\\* + &
\frac{\nu_{22} h}{J} \left(
  2 x_{\xi}x_{\eta} \pdiff{(r/h)}{\xi}
- 2 x_{\xi}x_{\xi} \pdiff{(r/h)}{\eta}
\right)
\Bigg]
\label{eq:eta_viscous_term}
\end{align}
After applying Green’s theorem with
$\vecn = (\nxi, \neta)^T$ the outward normal vector, and multiplied with $J$  the viscosity term of then $r$-momentum equation reads:
\begin{align}
= \Bigg[ &
\frac{\nu_{21}}{J} \left(
- 2 y_{\eta} y_{\eta} h \pdiff{(q/h)}{\xi}
+ 2 y_{\eta} y_{\xi} h \pdiff{(q/h)}{\eta}
\right)
+
\\
& + \frac{\nu_{22}}{J} \left(
- y_{\eta} y_{\eta} h \pdiff{(r/h)}{\xi}
+ y_{\eta} y_{\xi} h \pdiff{(r/h)}{\eta}
+ y_{\eta} x_{\eta} h \pdiff{(q/h)}{\xi}
- y_{\eta} x_{\xi} h \pdiff{(q/h)}{\eta}
\right)
+
\\
& + \frac{\nu_{21}}{J} \left(
+ x_{\xi} y_{\eta} h \pdiff{(r/h)}{\xi}
- x_{\xi} y_{\xi} h \pdiff{(r/h)}{\eta}
- x_{\xi} x_{\eta} h \pdiff{(q/h)}{\xi}
+ x_{\xi} x_{\xi} h \pdiff{(q/h)}{\eta}
\right)
+
\\
& + \frac{\nu_{22}}{J} \left(
- 2 x_{\eta} x_{\eta} h \pdiff{(q/h)}{\xi}
+ 2 x_{\eta} x_{\xi} h \pdiff{(q/h)}{\eta}
\right)
\Bigg] \nxi +
%---------------
\\
+ \Bigg[ &
\frac{\nu_{21}}{J} \left(
  2 y_{\xi} y_{\eta} h \pdiff{(q/h)}{\xi}
- 2 y_{\xi} y_{\xi} h \pdiff{(q/h)}{\eta}
\right)
\\* & +
\frac{\nu_{22}}{J} \left(
  y_{\xi}y_{\eta} h \pdiff{(r/h)}{\xi}
- y_{\xi}y_{\xi} h \pdiff{(r/h)}{\eta}
- y_{\xi}x_{\eta} h \pdiff{(q/h)}{\xi}
+ y_{\xi}x_{\xi} h \pdiff{(q/h)}{\eta}
\right)
\\* & +
\frac{\nu_{21}}{J} \left(
- x_{\xi}y_{\eta} h \pdiff{(r/h)}{\xi}
+ x_{\xi}y_{\xi} h \pdiff{(r/h)}{\eta}
+ x_{\xi}x_{\eta} h \pdiff{(q/h)}{\xi}
- x_{\xi}x_{\xi} h \pdiff{(q/h)}{\eta}
\right)
\\* & +
\frac{\nu_{22}}{J} \left(
2 x_{\xi}x_{\eta} h \pdiff{(r/h)}{\xi}
- 2 x_{\xi}x_{\xi} h \pdiff{(r/h)}{\eta}
\right)
\Bigg] \neta
\label{eq:eta_greens_viscous_term}
\end{align}
%
The linear functions in the first derivatives yields:
\begin{align}
    h \pdiff{(q/h)}{\xi} = \pdiff{q}{\xi} - \frac{q}{h}\pdiff{h}{\xi},
    & \qquad h \pdiff{q/h)}{\eta} = \pdiff{q}{\eta} - \frac{q}{h}\pdiff{h}{\eta},
    \\*
    h \pdiff{(r/h)}{\xi} = \pdiff{r}{\xi} - \frac{r}{h}\pdiff{h}{\xi},
    & \qquad h\pdiff{(r/h)}{\eta} = \pdiff{r}{\eta} - \frac{r}{h}\pdiff{h}{\eta}.
\end{align}
and will be used for the discretization of these viscosity terms.

%--------------------------------------------------------------------------------
 \subsubsection*{Cartesian grid}
On a cartesian grid ($y_\xi = x_\eta = 0$) these viscous terms reduces for the $q$-momentum equation (i.e.\ \autoref{eq:xi_greens_viscous_term}), to:
%
\begin{align}
    \Bigg[ &  \frac{\nu_{11}}{J} \left(
    - 2 y_{\eta} y_{\eta} h \pdiff{(q/h)}{\xi}
    \right)
    +
    \\* & +
    \frac{\nu_{12}}{J} \left(
    - y_{\eta}y_{\eta} h \pdiff{(r/h)}{\xi}
    - y_{\eta}x_{\xi} h \pdiff{(q/h)}{\eta} \right) +
    \Bigg] \nxi +
    %---------------
    \\
    + \Bigg[ &
    \frac{\nu_{11}}{J} \left(
    - x_{\xi}y_{\eta} h \pdiff{(r/h)}{\xi}
    - x_{\xi}x_{\xi} h \pdiff{(q/h)}{\eta}
    \right)
    \\* & +
    \frac{\nu_{12}}{J} \left(
    2 x_{\xi}x_{\eta} h \pdiff{(r/h)}{\xi}
    - 2 x_{\xi}x_{\xi} h \pdiff{(r/h)}{\eta}
    \right)
    \Bigg] \neta
\end{align}
%
And the $r$-momentum equation (i.e.\ \autoref{eq:eta_greens_viscous_term}) to:
%
\begin{align}
    = \Bigg[
    & + \frac{\nu_{22}}{J} \left(
    - y_{\eta} y_{\eta} h \pdiff{(r/h)}{\xi}
    - y_{\eta} x_{\xi} h \pdiff{(q/h)}{\eta}
    \right)
    +
    \\
    & + \frac{\nu_{21}}{J} \left(
    + x_{\xi} y_{\eta} h \pdiff{(r/h)}{\xi}
    + x_{\xi} x_{\xi} h \pdiff{(q/h)}{\eta}
    \right)
    \Bigg] \nxi +
    %---------------
    \\
    + \Bigg[ &
    \frac{\nu_{21}}{J} \left(
    - x_{\xi}y_{\eta} h \pdiff{(r/h)}{\xi}
    - x_{\xi}x_{\xi} h \pdiff{(q/h)}{\eta}
    \right)
    \\* & +
    \frac{\nu_{22}}{J} \left(
    - 2 x_{\xi}x_{\xi} h \pdiff{(r/h)}{\eta}
    \right)
    \Bigg] \neta
\end{align}
%
On a cartesian grid and $\nu_{12} = \nu_{21} = 0$:
\begin{align}
    \mat{\nu} =
    \begin{pmatrix}
        \nu_{11} & 0 \\
        0  & \nu_{22}
    \end{pmatrix}=
    \begin{pmatrix}
        \nu + \Psi_{11} & 0 \\
        0  & \nu + \Psi_{22}
    \end{pmatrix}
\end{align}
with $\Psi_{11}$ and $\Psi_{22}$ are the added artificial viscosity due to the regularization of the primairy variables (see \autoref{sec:artificial_viscosity}), for the $q$-equation yields
%
\begin{align}
    \Bigg[ &  \frac{\nu_{11}}{J} \left(
    - 2 y_{\eta} y_{\eta} h \pdiff{(q/h)}{\xi}
    \right)
    \Bigg] \nxi +
    %---------------
    \\*
    + \Bigg[ &
    \frac{\nu_{11}}{J} \left(
    - x_{\xi}y_{\eta} h \pdiff{(r/h)}{\xi}
    - x_{\xi}x_{\xi} h \pdiff{(q/h)}{\eta}
    \right)
    \Bigg] \neta
\end{align}
and for the $r$-momentum equation to:
\begin{align}
    = \Bigg[
    & \frac{\nu_{22}}{J} \left(
    - y_{\eta} y_{\eta} h \pdiff{(r/h)}{\xi}
    - y_{\eta} x_{\xi} h \pdiff{(q/h)}{\eta}
    \right)
    \Bigg] \nxi +
    %---------------
    \\
    + \Bigg[ &
    \frac{\nu_{22}}{J} \left(
    - 2 x_{\xi}x_{\xi} h \pdiff{(r/h)}{\eta}
    \right)
    \Bigg] \neta
\end{align}
with $\vecn = (\nxi, \neta)^T$ the outward normal vector.
%--------------------------------------------------------------------------------
\subsection{Discretization in time}
The discretization in time for the main terms is given and need to performed for each quadrature point at the eight faces of the control volume:
%
\begin{align}
    &\pdiff{q}{\xi} - \frac{q}{h}\pdiff{h}{\xi} \approx
    \\*
    & \approx
    \underbrace{ \left( \pdiff{q^{n+\theta,p}}{\xi} + \left(- \frac{q^{n+\theta,p}}{h^{n+\theta,p}}\right)   \pdiff{h^{n+\theta,p}}{\xi} \right)}_{\textit {to right hand side}} +
    %
    \\*
    & + \underbrace{ \frac{q^{n+\theta,p}}{(h^{n+\theta,p})^2} \pdiff{h^{n+\theta,p}}{\xi} }_{{\cal{A}}_q}  \theta \Delta h
    + \underbrace{\left(-\frac{1}{h^{n+\theta,p}} \pdiff{h^{n+\theta,p}}{\xi} \right)}_{{\cal{B}}_q}  \theta \Delta q
    \\*
    & +
    \underbrace{1}_{{\cal{C}}_q} \pdiff{}{\xi}(\theta \Delta q) +
    \underbrace{\left(-\frac{q^{n+\theta,p}}{h^{n+\theta,p}}\right)}_{{\cal{D}}_q}
    \pdiff{}{\xi}(\theta \Delta h)
\end{align}
%
\begin{align}
    &\pdiff{q}{\eta} - \frac{q}{h}\pdiff{h}{\eta}
    \\
& \approx
\underbrace{ \left( \pdiff{q^{n+\theta,p}}{\eta} + \left(- \frac{q^{n+\theta,p}}{h^{n+\theta,p}}\right)  \pdiff{h^{n+\theta,p}}{\eta} \right)}_{\textit {to right hand side}} +
%
\\*
& + \underbrace{ \frac{q^{n+\theta,p}}{(h^{n+\theta,p})^2} \pdiff{h^{n+\theta,p}}{\eta} }_{{\cal{A}}_q} \theta \Delta h
+ \underbrace{\left(-\frac{1}{h^{n+\theta,p}} \pdiff{h^{n+\theta,p}}{\eta}  \right)}_{{\cal{B}}_q} \theta \Delta q
\\*
& +
\underbrace{1}_{{\cal{C}}_q} \pdiff{}{\eta}(\theta \Delta q) +
\underbrace{\left(-\frac{q^{n+\theta,p}}{h^{n+\theta,p}}\right)}_{{\cal{D}}_q}
\pdiff{}{\eta}(\theta \Delta h)
\end{align}
%
\begin{align}
        &\pdiff{r}{\xi} - \frac{r}{h}\pdiff{h}{\xi} \approx
    \\*
& \approx
\underbrace{ \left( \pdiff{r^{n+\theta,p}}{\xi} + \left(- \frac{r^{n+\theta,p}}{h^{n+\theta,p}}\right)   \pdiff{h^{n+\theta,p}}{\xi} \right)}_{\textit {to right hand side}} +
%
\\*
& + \underbrace{ \frac{r^{n+\theta,p}}{(h^{n+\theta,p})^2} \pdiff{h^{n+\theta,p}}{\xi} }_{{\cal{A}}_r} \theta \Delta h
+ \underbrace{\left(-\frac{1}{h^{n+\theta,p}} \right) \pdiff{h^{n+\theta,p}}{\xi} }_{{\cal{B}}_r} \theta \Delta r
\\*
& +
\underbrace{1}_{{\cal{C}}_r} \pdiff{}{\xi}(\theta \Delta r) +
\underbrace{\left(-\frac{r^{n+\theta,p}}{h^{n+\theta,p}}\right)}_{{\cal{D}}_r}
\pdiff{}{\xi}(\theta \Delta h)
\end{align}
%
\begin{align}
    &\pdiff{r}{\eta} - \frac{r}{h}\pdiff{h}{\eta}\approx
    \\*
    & \approx
    \underbrace{ \left( \pdiff{r^{n+\theta,p}}{\eta} + \left(- \frac{r^{n+\theta,p}}{h^{n+\theta,p}}\right) \pdiff{h^{n+\theta,p}}{\eta} \right)}_{\textit {to right hand side}} +
    %
    \\*
    & + \underbrace{ \frac{r^{n+\theta,p}}{(h^{n+\theta,p})^2} \pdiff{h^{n+\theta,p}}{\eta} }_{{\cal{A}}_r}  \theta \Delta h
    + \underbrace{\left(-\frac{1}{h^{n+\theta,p}} \pdiff{h^{n+\theta,p}}{\eta} \right)}_{{\cal{B}}_r} \theta \Delta r
    \\*
    & +
    \underbrace{1}_{{\cal{C}}_r} \pdiff{}{\eta}(\theta \Delta r) +
    \underbrace{\left(-\frac{r^{n+\theta,p}}{h^{n+\theta,p}}\right)}_{{\cal{D}}_r}
    \pdiff{}{\eta}(\theta \Delta h)
\end{align}

%{{\color{red}
%        \begin{align}
%            \pdiff{}{x} =& \frac{y_{\eta}}{J}\pdiff{}{\xi}\left( \right)
%            - \frac{y_{\xi}}{J} \pdiff{}{\eta}\left( \right)
%            \\
%            \pdiff{}{y} =& \frac{-x_{\eta}}{J}\pdiff{}{\xi}\left( \right)
%            + \frac{x_{\xi}}{J} \pdiff{}{\eta}\left( \right)
%        \end{align}
%}}
